% Agent 05: Pages 361–364 (PDF pp. 9–11)
% Sections 2.3 (end), 2.4, and 2.5 (beginning)
% Equations (2.3.4)–(2.3.18), (2.4.1)–(2.4.9), (2.5.1)–(2.5.6)

%%% ====================================================================
%%% SECTION 2.3 (continued) — Cyclotron and Synchrotron Radiation
%%% ====================================================================

% [Continuation of Section 2.3]

\noindent
is
\begin{equation}
\frac{dE}{dt} = \frac{2\,r_0^2}{3\,c}\,(\gamma v_\perp)^2 B^2.
\tag{2.3.4}
\end{equation}
Here $v_\perp$ is the component of the electron velocity perpendicular to the magnetic
field.  If the electron is nonrelativistic, $v \ll c$, this radiation is called ``cyclotron
radiation''.  If the electron is ultrarelativistic, $\gamma \gg 1$, it is called ``synchrotron
radiation''.

\medskip
\noindent
\textit{Cyclotron radiation from a nonrelativistic plasma.}

\medskip
\noindent
Consider a plasma with temperature
\begin{equation}
kT \ll m_e c^2, \quad \text{i.e.}\ T \ll 6 \times 10^9\,\text{K},
\tag{2.3.5}
\end{equation}
and with a magnetic field of strength $B$. The radiation from protons spiralling
in the magnetic field is smaller by a factor $(m_e/m_p)^3 \simeq 10^{-10}$ than that from
electrons, since
\begin{equation}
\frac{dE}{dt} \propto \frac{v^2}{m^3} \propto \frac{1}{m}.
\tag{2.3.6}
\end{equation}
(Recall: in thermal equilibrium the mean kinetic energies of protons and electrons
are the same.)  Hence, we can ignore protons and other ions when calculating
cyclotron radiation.  Let $f_e$ be the number of unbound electrons per baryon in
the plasma.  Then the total emissivity (ergs per second per gram of plasma) is
\begin{equation}
\varepsilon = f_e \left\langle \frac{dE}{dt} \right\rangle \bigg/ m_p,
\tag{2.3.7}
\end{equation}
where $\langle\;\rangle$ denotes an average over the Maxwell--Boltzmann velocity distribution
of the electrons.  Since $\gamma = 1$, since 2 of the 3 spatial directions are orthogonal
to $B$, and since the mean kinetic energy per electron is $\frac{3}{2}kT$, we have
\begin{equation}
\langle (\gamma v_\perp)^2 \rangle = \frac{2}{3}\langle v^2 \rangle = \frac{2}{3}\cdot\frac{2kT}{m_e}.
\tag{2.3.8}
\end{equation}
Combining this with equations~(2.3.4) and~(2.3.7), we obtain
\begin{align}
\varepsilon &= \frac{4}{3}\left(\frac{r_0^2 c}{m_p}\right)
\left(\frac{kT}{m_e c^2}\right)B^2 \notag \\
&= (0.32~\text{ergs/g\,sec})\,f_e\,T_K\,B_G^2,
\tag{2.3.9}
\end{align}
where $T_K$ is temperature in ${}^\circ$K and $B_G$ is magnetic field in Gauss.

Each electron emits its radiation monochromatically, with the cyclotron
frequency (orbital frequency of electron's spiral motion)
\begin{equation}
\nu_{\text{cyc}} = \frac{eB}{2\pi m_e c} = (2.79~\text{MHz})\,T_K\,B_G.
\tag{2.3.10}
\end{equation}

Thus, if the magnetic field is uniform, the radiation will be monochromatic
with this frequency.  But in Nature the magnetic field will always be inhomo\-
geneous, and the spectrum will show a spread proportional to the spread in $B^2$.
(Of course, there are always other sources of spread in the spectrum---e.g., doppler
shifts and relativistic generation of harmonics.)

\medskip
\noindent
\textit{Synchrotron radiation from a relativistic plasma.}

\medskip
\noindent
Consider a plasma with temperature
\begin{equation}
kT \gg m_e c^2, \quad \text{i.e.}\ T \gg 6 \times 10^9\,\text{K},
\tag{2.3.11}
\end{equation}
and with a magnetic field of strength $B$.
As in the nonrelativistic case, radiation
from protons and ions is totally negligible.
The total emissivity due to electrons
and positrons (recall: at $kT \gg m_e c^2$, there will be many electron-positron pairs!)
is given, as before, by equation~(2.3.7); but now the factor $f_e$ must be the number
of electrons and positrons per baryon,\footnote{The case of an optically thin plasma
is of particular interest in astrophysics. The kinetic theory of the creation and
annihilation of positrons in this case is somewhat complex; see, e.g.,
\citet{BisnovatyiKogan1971}.}
and the relevant, ultrarelativistic Maxwell--Boltzmann average is

\begin{equation}
\langle (\gamma v_\perp)^2 \rangle = \tfrac{2}{3}\langle\gamma^2\rangle
= 8\left(\frac{kT}{m_e c^2}\right)^{\!2}.
\tag{2.3.12}
\end{equation}
[One evaluates $\langle\gamma^2\rangle$ using the distribution function in phase space
\begin{equation}
\mathscr{N} = \frac{dN}{d^3 x\,d^3 p} \propto e^{-E/kT} = e^{-\gamma m_e c^2/kT},
\tag*{}
\end{equation}
with
\begin{equation}
d^3 p = 4\pi c^{-3}E^2\,dE
\tag*{}
\end{equation}
in ultrarelativistic limit;
in particular
\begin{equation}
\langle\gamma^2\rangle = \frac{1}{(m_e c^2)^2}\,
\frac{\displaystyle\int_0^\infty E^2\,e^{-E/kT}\,E^2\,dE}
{\displaystyle\int_0^\infty e^{-E/kT}\,E^2\,dE}
= 12\left(\frac{kT}{m_e c^2}\right)^{\!2}.\,]
\tag*{}
\end{equation}

Combining equations~(2.3.12), (2.3.7), and~(2.3.4), we obtain for the total
emissivity of the plasma
\begin{align}
\varepsilon &= \frac{16}{3}\left(\frac{r_0^2 c}{m_p}\right)
\left(\frac{kT}{m_e c^2}\right)^{\!2} B^2 \notag \\
&= (2.2 \times 10^{-10}~\text{ergs/g\,sec})\,f_e\,T_K^2\,B_G^2.
\tag{2.3.13}
\end{align}

In the ultrarelativistic case the electrons do not emit monochromatically.
An electron of energy $\gamma m_e c^2$ beams most of its radiation into a forward cone of
half-angle $\alpha = 1/\gamma$ (special-relativistic ``headlight effect'').
Since the electron is
headed toward the observer with speed $v$ when its cone is directed toward him,
the cone sweeps past the observer in time
\begin{equation}
\Delta t = (1-v^2)\left(\frac{2\alpha}{\omega_{\text{cyc}}}\right) = (1-v^2)\,\frac{1}{\gamma^3\,\omega_{\text{cyc}}}
= \frac{1}{\gamma^3\,\omega_{\text{cyc}}}.
\tag{2.3.14}
\end{equation}
Here $\omega_{\text{cyc}}$ is the angular velocity of the electron in its spiralling orbit
\begin{equation}
\omega_{\text{cyc}} = \frac{eB}{\gamma m_e c}.
\tag{2.3.15}
\end{equation}

Thus, the radiation comes in short bursts, of duration $\sim\!\Delta t = 1/(\gamma^3\omega_{\text{cyc}})$,
separated by long intervals of time $2\pi/\omega_{\text{cyc}}$.
When Fourier analyzed, such
radiation must be concentrated near the ``critical frequency''
\begin{equation}
\nu_{\text{crit}} = \tfrac{3}{2}\,\nu_{\text{cyc}} = \gamma^2\,eB/m_e c.
\tag{2.3.16}
\end{equation}

A detailed calculation [chapter~14 of \citet{Jackson1962}] reveals a fairly broad
spectrum which rises as $\nu^{1/3}$ at low frequencies, which peaks at
\begin{equation}
\nu_{\text{peak}} = 0.29\,\nu_{\text{crit}},
\tag{2.3.17}
\end{equation}
and which decays exponentially $\sim\nu^{1/2}\,e^{-2\nu/\nu_{\text{crit}}}$, at $\nu \gg \nu_{\text{crit}}$.
When averaged over the relativistic plasma, such a spectrum will have a broad
peak, with maximum located at
\begin{equation}
\nu_{\text{peak, plasma}} = \frac{\gamma_{\text{peak}}^2\,eB}{4\,m_e c}
\simeq 12(kT/m_e c^2)^2
\quad\left[\text{value of $\gamma_{\text{peak}}^2$ for electrons of}\right]
\tag{2.3.18}
\end{equation}
\begin{equation}
\nu_{\text{peak, plasma}} \simeq (100~\text{MHz})\,T_{10}^2 B_G.
\tag*{}
\end{equation}
Here $T_{10}$ is temperature in units of $10^{10}$\,K.
At $\nu \ll \nu_{\text{peak, plasma}}$
the spectrum will rise as a power law.

%%% ====================================================================
%%% SECTION 2.4 — Electron Scattering of Radiation
%%% ====================================================================

\subsection{Electron Scattering of Radiation}
\label{sec:2.4}

\textit{Basic references}

\citet{Jackson1962}, \S\,14.7 of \citet{Jackson1962b};
\S\,12.3 of \citet{Leighton1959};
\citet{Kompaneets1957}; \citet{Weyman1965};
\citet{Sunyaev1973}.

\medskip
\noindent
\textit{Basic physics and formulas}

\medskip
\noindent
In discussing electron scattering, we shall confine attention to the nonrelativistic
case---i.e.\ we shall demand that the temperature of the gas and the photon
frequencies satisfy
\begin{equation}
kT \ll m_e c^2 \quad (T \ll 6 \times 10^9\,\text{K}); \quad
h\nu \ll m_e c^2 = 500~\text{keV}.
\tag{2.4.1}
\end{equation}

The differential cross section for a free electron to scatter a photon has the
form
\begin{equation}
\frac{d\sigma_e}{d\Omega} = \tfrac{1}{2}\,r_0^2\,(1 + \cos^2\theta),
\tag{2.4.2}
\end{equation}
where $\theta$ is the angle between incoming and outgoing photon.  The total
cross section (``Thomson cross section''), obtained by integrating over all
directions, is
\begin{equation}
\sigma_T = (8\pi/3)\,r_0^2 = 0.665 \times 10^{-24}~\text{cm}^2.
\tag{2.4.3}
\end{equation}
It is important for astrophysical applications that the scattering cross section
is ``color blind'' (no dependence on frequency).  See, e.g., \S\,4.

It is also important that the differential cross section is symmetric between
forward and backward directions.  This guarantees, for example, that on the
average a scattered photon transmits \textit{all} of its momentum to the electron
(Here $\mathbf{n}_s$ is a unit vector in the direction of the initial photon motion).  By
contrast, only a tiny fraction of the photon energy is transmitted to the electron.
In a frame where the electron is initially at rest, it receives a kinetic energy
\begin{subequations}
\begin{equation}
\Delta E_\gamma = (h\nu/m_e c^2)\,h\nu\,(1-\cos\theta).
\tag{2.4.4a}
\end{equation}
\begin{equation}
(\Delta\mathbf{p}_e) = \mathbf{p}_\gamma = (h\nu/c)\,\mathbf{n}_s.
\tag{2.4.4b}
\end{equation}
\end{subequations}
(Recall: we have assumed $h\nu/m_e c^2 \ll 1$.)

Consider a monochromatic beam of photons, with frequency $\nu$, moving
through a thermalized plasma of temperature $T$.  Scattering by protons and ions
will be negligible compared to scattering by electrons, since
\begin{equation}
\sigma_e \propto r_0^2 = (e^2/m_e c^2)^2 \propto 1/(\text{mass of scatterer})^2.
\tag{2.4.5}
\end{equation}

On the average, how much energy $\langle\Delta E_\gamma\rangle$ is transmitted to the electron in a
single scattering?  The answer for electrons nearly at rest ($kT \ll h\nu$) is obtained
by averaging expression~(2.4.4b) over all angles.  Because of the forward-backward
symmetry of the differential cross section, $\cos\theta$ averages to zero and one obtains
\begin{equation}
\langle\Delta E_\gamma\rangle = (h\nu/m_e c^2)\,h\nu, \quad \text{if } kT \ll h\nu.
\tag{2.4.6}
\end{equation}

One can calculate $\langle\Delta E_\gamma\rangle$ for higher temperatures by invoking conservation
of 4-momentum, and averaging over all angles and electron speeds.  However,
such a calculation is rather long and messy.  To obtain the answer more easily,
one can use the following trick: An examination of the law of 4-momentum
conservation convinces one that $\langle\Delta E_\gamma\rangle$ must have the temperature dependence
\begin{equation}
\langle\Delta E_\gamma\rangle \propto (h\nu - \alpha kT),
\tag{2.4.7}
\end{equation}
where $\alpha$ is a constant to be calculated.  (Thus, if $h\nu > \alpha kT$, the electrons get
heated by the photons; if $h\nu < \alpha kT$, they get cooled.)
This law will reduce to
expression~(2.4.6) for $kT \ll h\nu$ if and only if the proportionality constant is
$h\nu/m_e c^2$:
\begin{equation}
\langle\Delta E_\gamma\rangle = (h\nu/m_e c^2)(h\nu - \alpha kT).
\tag*{}
\end{equation}

To calculate the constant $\alpha$, imagine the following experiment.  Place a large
number of photons and electrons into a box with perfectly reflective walls.
Require that the photons and electrons interact only by electron scattering.  Then
photons cannot be created or destroyed; only scattered.  Hence, when thermal
equilibrium is reached, the photons acquire a Boltzmann energy distribution,
rather than a Planck distribution.  (Recall that stimulated emissions are responsible
for Planckian deviations from the Boltzmann law; \S\,2.5.)  Since photons have
zero rest mass, their Boltzmann distribution
\begin{equation}
\mathscr{N} = \frac{dN}{d^3 x\,d^3 p} \propto e^{-h\nu/kT},
\tag*{}
\end{equation}
corresponds to a number of photons per unit frequency given by
\begin{equation}
\frac{dN}{d\nu} \propto \nu^2\,e^{-h\nu/kT},
\tag{2.4.8}
\end{equation}
and corresponds to
\begin{equation}
\langle h\nu\rangle = 3kT, \qquad
\langle (h\nu)^2\rangle = 12(kT)^2.
\tag*{}
\end{equation}
Thus, for our thought experiment the energy transfer in each collision, expression
(2.4.7) averaged over the equilibrium photon distribution, is
\begin{equation}
\langle\langle\Delta E_\gamma\rangle\rangle = (3kT/m_e c^2)(4-\alpha)\,kT.
\tag*{}
\end{equation}

But in equilibrium there must be no average energy transfer between photons
and electrons; $\langle\langle\Delta E_\gamma\rangle\rangle$ must be zero.
This condition tells us that $\alpha = 4$.
Thus, turning back to the general situation, we conclude that monochromatic
photons passing through a plasma of temperature $T$ transfer an average energy
per collision
\begin{equation}
\langle\Delta E_\gamma\rangle = (h\nu/m_e c^2)(h\nu - 4kT)
\tag{2.4.9}
\end{equation}
to the electrons.  When $4kT \gg h\nu$, the photon energies get boosted by the
collisions, and one says that the radiation is being ``Comptonized''; see \S\,5.10;
see, e.g., \citet{Illarionov1972}.

This concludes, for the moment, our discussion of the interaction between
radiation and plasmas.  We must point out that we have ignored a number of
plasma effects and instabilities which can be important in astrophysical
situations.  See, e.g., \citet{Bekefi1966} and \citet{KaplanTsytovich1972}.

%%% ====================================================================
%%% SECTION 2.5 — Hydrodynamics and Thermodynamics
%%% ====================================================================

\subsection{Hydrodynamics and Thermodynamics}
\label{sec:2.5}

\textit{Basic references}\quad
\S\S\,22.2 and 22.3 of \citet{MTW1973};
\citet{Lichnerowicz1967} and references cited therein.
We adopt the notational conventions of \citet{MTW1973},
including signature ``$-+++$''; geometrized units ($c = G = 1$); Greek
indices ranging from 0 to 3 and Latin from 1 to 3; ``hats'' on indices that refer
to local orthonormal frames, e.g.\ $\mu^{\hat{\alpha}}$ and $T^{\hat{\alpha}\hat{\beta}}$; extra-bold, sans-serif type for
4-vectors and 4-tensors, e.g.\ $\fvec{u}$ and $\mathbf{T}$; normal bold-face type for 3-vectors and
3-tensors (local Euclidean geometry).

\medskip
\noindent
\textit{Parameters describing the ``fluid.''}

\medskip
\noindent
``LRF'' $=$ local rest frame of the baryons; i.e.\ local orthonormal frame in which
there is no net baryon flux in any direction.

\begin{itemize}
\item[$\fvec{u}$] 4-velocity of LRF; i.e.\ 4-velocity of the ``fluid''.
\item[$n$] number density of baryons (number of baryons per unit volume) as
measured in LRF.
\item[$m_B$] mean rest mass of a baryon in the fluid; i.e.
\end{itemize}
\begin{equation}
m_B = \left(\!\!
\begin{array}{c}
\text{rest mass of}\\
\text{hydrogen atom}
\end{array}\!\!\right)
\times
\left(\!\!
\begin{array}{c}
\text{fraction of all baryons that are}\\
\text{in the form of hydrogen nuclei}
\end{array}\!\!\right)
+
\frac{1}{4}
\left(\!\!
\begin{array}{c}
\text{rest mass of}\\
\text{helium atom}
\end{array}\!\!\right)
\times
\left(\!\!
\begin{array}{c}
\text{fraction of all baryons that}\\
\text{are in helium nuclei}
\end{array}\!\!\right)
+ \cdots
\tag{2.5.1}
\end{equation}

\textit{Note}: in this section we restrict ourselves to fluids that are chemically
homogeneous and in which no nuclear reactions occur (``standard
fluid'').  Thus, $m_B$ is constant.

\begin{align}
\rho_0 &\equiv m_B\,n, && \text{rest-mass density, defined by}
\tag{2.5.2} \\[4pt]
\mathcal{V} &\equiv 1/n, && \text{specific volume, i.e.\ volume per baryon, defined by}
\tag{2.5.3} \\[4pt]
\mathcal{V}_0 &= 1/\rho_0, && \text{specific volume, i.e.\ volume per unit rest mass, defined by}
\tag{2.5.4}
\end{align}

\begin{align}
\rho &= \rho_0(1 + \Pi), && \text{total density of mass-energy, as measured in LRF.}
\tag{2.5.5}\\[4pt]
\Pi &\quad && \text{specific internal energy, defined by}
\tag*{}
\end{align}

\textit{Note}: here and throughout this section we set the speed of light equal to
one.
\begin{align}
p  &\quad && \text{isotropic pressure, as measured in LRF.}
\tag*{} \\
T  &\quad && \text{temperature, as measured in LRF.}
\tag*{} \\
s  &\quad && \text{entropy per baryon, as measured in LRF; of course,}
\tag*{} \\
s_0 &= m_B^{-1}\,s. && \text{entropy per unit mass, as measured in LRF.}
\tag{2.5.6}
\end{align}
